% SPDX-License-Identifier: CC-BY-4.0
% Copyright 2025 Florian Aram Feuerriegel — kassensturz.org
% Abschnitt 2 — Das Modell (js/rechner/berechne.js, einkommensteuer.js,
% verteilung.js, transition.js, abgeleitet.js; Daten in js/data.js)

\section{Modell}
\label{sec:modell}

\subsection{Aufbau}
\label{sec:aufbau}

Das Modell verbindet eine statische Mikrosimulation mit einem dynamischen
Makrorahmen. Zwölf Haushaltstypen $i = 1,\dots,12$ stehen für die
Einkommensdezile D1 bis D9 und drei Teile des obersten Dezils (D10a: 90.
bis 95.~Perzentil, D10b: 95. bis 99., D10c: oberstes Prozent); $N_i$ ist die
Zahl der Haushalte des Typs in Millionen, zusammen rund 41~Millionen. Die Zeit
läuft in Perioden $t = 1,\dots,5$ zu je $n = 4$ Jahren von 2025 bis 2044;
$j_t$ ist das erste Jahr der Periode. Innerhalb einer Periode gilt ein
Politikvektor $\theta$ unverändert. Der Status quo $\theta^0$ ist der
Rechtsstand 2026 (Tabelle~\ref{tab:stellgroessen}).

Jede Periode durchläuft die Schritte in Abbildung~\ref{fig:ablauf}. Aus dem
Politikvektor und dem Zustand zu Periodenbeginn entstehen Einkommen, Steuern
und Transfers je Haushaltstyp, daraus Aufkommen, Ausgaben und Saldo des
Staates, die Nachfragewirkung und die Verteilungsmaße. Der Übergang schreibt
BIP, Schuldenquote, kumulierte Emissionen und die Kapitalniveaus in die
nächste Periode fort. Ergebnisse sind Differenzen zum Status quo derselben
Periode, $\Delta x_t = x_t(\theta) - x_t(\theta^0)$.

Geldgrößen des Staates stehen in Mrd.~\euro{} je Jahr, Geldgrößen der
Haushalte in \euro{} je Haushalt und Jahr, beide in Preisen von 2025, soweit
nicht anders angegeben. Sätze sind Anteile; Änderungen werden in
Prozentpunkten (Pp.) berichtet. Summen $\sum_i N_i x_i$ über Größen $x_i$ in
Euro je Haushalt sind in Mrd.~\euro{} angegeben; der Faktor $10^{-3}$ wird
nicht geschrieben.

\begin{figure}[t]
\centering
\begin{tikzpicture}[
  node distance=3.5mm and 4mm,
  box/.style={draw=grau, rounded corners=2pt, fill=hell, align=center, font=\footnotesize,
              minimum height=8.5mm, text width=26mm},
  pfeil/.style={-{Stealth[length=1.8mm]}, semithick, grau}]
\node[box] (a) {1 Geltendes Recht\\einrechnen};
\node[box, right=of a] (b) {2 Arbeitsangebot};
\node[box, right=of b] (c) {3 Einkommensteuer};
\node[box, right=of c] (d) {4 Weitere Steuern};
\node[box, below=of d] (e) {5 Sozialbeiträge};
\node[box, left=of e] (f) {6 Transfers\\und Renten};
\node[box, left=of f] (g) {7 Haushalte:\\Netto, Verteilung};
\node[box, left=of g] (h) {8 Nachfrage\\der Periode};
\node[box, below=of h] (i) {9 Fortschreibung\\in Euro der Periode};
\node[box, right=of i] (j) {10 Saldo,\\Schuldenbremse};
\node[box, right=of j, fill=white, text width=56mm] (k) {Übergang zur nächsten Periode:\\BIP, Schulden, Emissionen, Kapital};
\draw[pfeil] (a)--(b); \draw[pfeil] (b)--(c); \draw[pfeil] (c)--(d); \draw[pfeil] (d)--(e);
\draw[pfeil] (e)--(f); \draw[pfeil] (f)--(g); \draw[pfeil] (g)--(h); \draw[pfeil] (h)--(i);
\draw[pfeil] (i)--(j); \draw[pfeil] (j)--(k);
\end{tikzpicture}
\caption{Ablauf einer Periode (\datei{berechne()}) und Übergang (\datei{berechneTransition()}).}
\label{fig:ablauf}
\end{figure}

\paragraph{Geltendes Recht.} Beschlossene Änderungen, die über die Jahre
wirken, kommen zu den Stellgrößen hinzu, gemittelt über die Jahre $j$ der
Periode. Die Körperschaftsteuer sinkt ab 2028 um einen Prozentpunkt je Jahr
bis auf 10\,\% (§\,23 Abs.~1 KStG), der Beitragssatz der Rentenversicherung
folgt nach §\,158 SGB~VI den Ausgaben, 18,6\,\% bis 2027, 20,1\,\% im Jahr 2030
und 21,15\,\% ab 2039 \parencite{bmas2025}:
\begin{gather}
  \tau\hoch{KSt}_t = \max\Bigl\{0,\ \tau\hoch{KSt} - \frac1n\sum_{j\in t}\min\bigl\{0{,}05,\ \max\{0,\ (j-2027)/100\}\bigr\}\Bigr\},\\
  \beta\hoch{RV}_t = \beta\hoch{RV} + \frac1n\sum_{j\in t}\bigl(\bar\beta\hoch{RV}(j) - 0{,}186\bigr),
\end{gather}
mit $\bar\beta\hoch{RV}(j)$ linear zwischen den genannten Stützstellen. Die
Stellgrößen zeigen die Werte von 2026.

\subsection{Haushalte}
\label{sec:haushalte}

Ein Haushaltstyp ist beschrieben durch das Bruttoeinkommen $B_i$, den
Kapitalanteil $\kappa_i$, die durchschnittliche Konsumquote $c_i$, das
Bedarfsgewicht $w_i$ der neuen OECD-Skala, das Vermögen $V_i$ und den Anteil
$\rho_i$ der gesetzlichen Rente am Einkommen (Tabelle~\ref{tab:dezile} im
Anhang). Grundlage sind SOEP~v40, der Mikrozensus 2024, die
Verteilungsrechnung DINA-DE und der Vermögensbericht des DIW. Arbeits- und
Kapitaleinkommen sind $A_i = (1-\kappa_i)B_i$ und $K_i = \kappa_i B_i$. Das
oberste Dezil ist dreigeteilt, weil sich dort Spitzensteuersatz und
Vermögensteuern konzentrieren: D10c hat im Mittel 700.000\,\euro{}
Bruttoeinkommen, davon 45\,\% aus Kapital, und 7~Mio.\,\euro{} Vermögen.

Haushalte und Staat benutzen dieselben Profile: Kindergeldkinder $k_i$
(zusammen 17~Mio.), Bürgergeld-Regelsatzjahre $q_i$ (zusammen 3,9~Mio., nur
D1 bis D4), die Grenzkonsumneigung $m_i$ von 0,65 in D1 bis 0,15 in D10c
\parencite{jappelli2014,fagereng2021} und das Gewicht $\omega_i$ der
\COz-Last mit $\sum_i N_i\omega_i = 1$. Es verteilt die Last regressiv: Bei
55\,\euro/t trägt D1 \LastDeins\,\% seines Bruttoeinkommens, D10c \LastDzehnc\,\%.

\subsection{Arbeitsangebot}
\label{sec:arbeit}

Das Einkommen reagiert auf den Anteil $1-\mu_i$, der vom letzten verdienten
Euro bleibt:
\begin{equation}
  \lambda_i = \max\Bigl\{0{,}55,\ a_i\cdot\min\Bigl\{1{,}25,\ \max\Bigl\{0{,}55,\ 1 + \varepsilon_i\,\frac{\mu_i^0 - \mu_i}{1 - \mu_i^0}\Bigr\}\Bigr\}\Bigr\},
  \qquad \tilde A_i = \lambda_i A_i .
\end{equation}
Dabei ist $\mu_i$ der Grenzsteuersatz des Haushalts (Abschnitt~\ref{sec:est})
unter $\theta$ und $\mu_i^0$ im Status quo, beide am Arbeitseinkommen $A_i$
ausgewertet. Die Elastizität $\varepsilon_i = 0{,}20$ entspricht dem Konsens
zur intensiven Marge \parencite{gruber2002,saez2012}. Für D10c gilt
$\varepsilon_{12} = 0{,}40$, weil Spitzeneinkommen vor allem durch Verlagerung
und Verhandlung reagieren \parencite{piketty2014,kleven2014}. Dazu kommt ein
Ausweichfaktor, den der tarifliche Spitzensatz $r_5$ auslöst: Steuervermeidung
oberhalb von 45\,\%, Wegzug oberhalb von 60\,\%, beides nur für den Anteil
$\varsigma = \min\{1,\ (x_m/x^S)^{\alpha-1}\}$ des Einkommens, der im Pareto-Rand
über der Grenze der Spitzenzone liegt (Abschnitt~\ref{sec:est}; im Status quo rund
die Hälfte):
\begin{equation}
  a_{12} = \max\bigl\{0{,}40,\ 1 - \varsigma\,\bigl(0{,}5\max\{0,\ r_5-0{,}45\} + 0{,}1\max\{0,\ r_5-0{,}60\}\bigr)\bigr\};
\end{equation}
für alle übrigen Typen ist $a_i = 1$. Die Reaktion betrifft nur die
Arbeitseinkünfte; Kapitaleinkünfte bleiben wie in den Daten, weil für sie keine
belegte Elastizität vorliegt, und $\tilde B_i = \tilde A_i + K_i$.

\subsection{Einkommensteuer}
\label{sec:est}

Der Tarif ist der Formeltarif des §\,32a EStG 2026, geschrieben als Integral
einer stückweise linearen Grenzsteuerrate $r(x)$ über das zu versteuernde
Einkommen $x$. Mit dem Grundfreibetrag $F$, dem Beginn der Spitzenzone $x^S$,
dem Eingangssatz $r_0$, dem Spitzensatz $r_5$ und $\xi = x - F$ gilt
\begin{equation}
  r(x) = \begin{cases}
    0 & \xi \le 0,\\
    r_0 + (r_m - r_0)\,\xi/b_2 & 0 < \xi \le e_2,\\
    r_m + (r_4 - r_m)\,(\xi - e_2)/b_3 & e_2 < \xi \le e_3,\\
    r_4 & e_3 < \xi \le x^S - F,\\
    r_5 & \xi > x^S - F,
  \end{cases}
  \qquad T(x) = \int_0^x r(z)\,dz,
\end{equation}
mit den Zonenbreiten von 2026, $b_2 = 5.451$\,\euro{} und $b_3 = 52.079$\,\euro,
den Zonenenden $e_2 = \min\{b_2,\ x^S - F\}$ und
$e_3 = \min\{b_2 + b_3,\ x^S - F\}$, $r_4 = \min\{0{,}42,\ r_5\}$ und
$r_m = r_0 + 0{,}35607\,(r_4 - r_0)$. Im Status quo ($F = 12.348$\,\euro,
$r_0 = 14\,\%$, $r_5 = 45\,\%$, $x^S = 277.826$\,\euro) weicht $T$ an keiner
Stelle um mehr als 1\,\% vom amtlichen Tarif ab. Jede Stellgröße wirkt nur auf
ihre Zone: $F$ verschiebt die Zonen~2 bis~4, $r_0$ den Beginn der Progression,
$r_5$ die Zone oberhalb von $x^S$ und $x^S$ nur deren Beginn. Liegt $r_5$
unter 42\,\%, wird die Proportionalzone auf $r_5$ gedeckelt.

Die Daten sind Haushaltseinkommen, der Tarif gilt je Steuerpflichtigem. Mit
dem Anteil $\zeta = 0{,}79$ des zu versteuernden Einkommens am Brutto
(Werbungskosten, Vorsorgeaufwand) und $\eta = 1{,}6$ Tarifeinheiten je
Haushalt (Einzelpersonen und zusammenveranlagte Paare) ist
\begin{equation}
  T_i = \eta\,T\bigl(\zeta\tilde A_i/\eta\bigr) + \tau\hoch{Abg}K_i,
  \qquad \mu_i = \zeta\,r\bigl(\zeta A_i/\eta\bigr),
\end{equation}
mit der Abgeltungsteuer $\tau\hoch{Abg}$ auf Kapitaleinkommen. Die Steuer
greift auf das Einkommen nach, der Grenzsteuersatz auf das Einkommen vor der
Reaktion des Arbeitsangebots zu. Im obersten Prozent liegt das mittlere zu
versteuernde Einkommen, rund 190.000\,\euro, unter $x^S$; mit dem
Durchschnitt allein zahlte niemand den Spitzensatz. Die Einkommen in D10c
folgen deshalb einer Pareto-Verteilung mit dem Exponenten $\alpha = 1{,}5$,
dem schweren Rand der Schätzungen für Deutschland
\parencite{atkinson2011,bartels2015}. Mit dem Mittelwert $\bar x$ der
Verteilung, ihrem Skalenparameter $x_m = \bar x(\alpha-1)/\alpha$, der unter
$x^S$ liegt, und dem Tarif ohne Spitzenzone $T_{\le 4}$ gilt
\begin{equation}
  T_{12} = \eta\Bigl[T_{\le 4}(\bar x) + (r_5 - r_4)\,\frac{x_m^{\alpha}\,(x^S)^{1-\alpha}}{\alpha - 1}\Bigr] + \tau\hoch{Abg}K_{12},
  \qquad
  \mu_{12} = \zeta\Bigl[r_{\le 4}(\bar x) + (r_5 - r_4)\Bigl(\frac{x_m}{x^S}\Bigr)^{\alpha}\Bigr];
\end{equation}
der Bruch ist der erwartete Einkommensteil oberhalb von $x^S$, der letzte
Faktor der Anteil der Veranlagungen dort. Das Aufkommen ist
$R\hoch{ESt} = \sum_i N_i T_i$.

\subsection{Weitere Steuern und Beiträge}
\label{sec:steuern}

\paragraph{Umsatzsteuer.} Der Konsum eines Haushalts ist
$C_i = c_i\,(\tilde B_i - T_i - \tfrac12\sigma_i)$ mit dem Sozialbeitrag
$\sigma_i$ (unten). 70\,\% davon tragen den Regelsatz $\tau\hoch{USt}$,
30\,\% den ermäßigten Satz $\tau\hoch{erm}$. Der besteuerte Konsum reagiert
mit der Semi-Elastizität $\varepsilon_C = -0{,}35$ je Einheit des Satzes:
\begin{equation}
\begin{split}
  R\hoch{USt} &= \varphi\,(1-\gamma)\sum_i N_i C_i
  \Bigl[0{,}7\,\frac{\tau\hoch{USt}}{1+\tau\hoch{USt}}\,\chi\bigl(\Delta\tau\hoch{USt}\bigr)
      + 0{,}3\,\frac{\tau\hoch{erm}}{1+\tau\hoch{erm}}\,\chi\bigl(\Delta\tau\hoch{erm}\bigr)\Bigr],\\
  \chi(\Delta) &= \min\{1{,}1,\ \max\{0{,}85,\ 1 + \varepsilon_C\Delta\}\}.
\end{split}
\end{equation}
Der Konsum der Haushalte erfasst nur rund 60\,\% der Bemessungsgrundlage; der
Faktor $\varphi = 1{,}68$ ergänzt die Käufe von Staat, Wohnungsbau und
befreiten Branchen, $\gamma = 3{,}7\,\%$ ist die Umsatzsteuerlücke.

\paragraph{Unternehmensteuern.} Körperschaft- und Gewerbesteuer haben
getrennte Bemessungsgrundlagen, $B\hoch{KSt} = 300$ und
$B\hoch{GewSt} = 536$~Mrd.~\euro{}, weil auch Personengesellschaften
Gewerbesteuer zahlen. Beide schrumpfen mit dem Investitionsfaktor
\begin{equation}
\begin{split}
  I &= 1 + \varepsilon_I\bigl(\tau\hoch{KSt} + \tau\hoch{GewSt} - 0{,}30\bigr),
  \qquad \bar I = \min\{1{,}2,\ \max\{0{,}7,\ I\}\},\\
  R\hoch{KSt} &= B\hoch{KSt}\,\bar I\,\tau\hoch{KSt},
  \qquad
  R\hoch{GewSt} = B\hoch{GewSt}\,\bar I\,\tau\hoch{GewSt},
\end{split}
\end{equation}
mit $\varepsilon_I = -0{,}40$: Zehn Prozentpunkte Unternehmensteuerlast über
30\,\% senken Gewinne und Investitionen um 4\,\%. Die Meta-Analyse von
\textcite{gechert2022} findet Wachstumswirkungen, die kleiner sind als oft
behauptet. Der Faktor wirkt sofort auf die Bemessungsgrundlage und
langfristig auf das private Kapital (Abschnitt~\ref{sec:dynamik}).

\paragraph{Vermögensbezogene Steuern.} Die Erbschaftsteuer rechnet mit einer
Erbmasse $H = 400$~Mrd.~\euro{} je Jahr, von der nach den persönlichen
Freibeträgen der Anteil $h = 0{,}45$ steuerpflichtig bleibt. 60\,\% der Masse
tragen den effektiven Satz $\tau\hoch{ErbSt}_{\mathrm{eff}}$, bei Verschonung
des Betriebsvermögens (Status quo) $0{,}3\,\tau\hoch{ErbSt}$, sonst
$0{,}9\,\tau\hoch{ErbSt}$; die übrigen 40\,\% tragen den halben, auf 15\,\%
begrenzten Satz. Die Vermögensteuer trifft Vermögen über
2~Mio.\,\euro{} ($W = 3.500$~Mrd.~\euro), deren deklarierte Höhe mit dem Satz
sinkt \parencite{bruelhart2022}, die Bodenwertsteuer den Bodenwert
($L = 5.000$~Mrd.~\euro). Die Mindeststeuer auf Milliardärsvermögen
($W^{*} = 600$~Mrd.~\euro) rechnet die gezahlte Einkommensteuer von rund
0,3\,\% des Vermögens und eine Vermögensteuer an \parencite{zucman2024}:
\begin{align}
  R\hoch{ErbSt} &= h\,H\bigl[0{,}6\,\tau\hoch{ErbSt}_{\mathrm{eff}} + 0{,}2\min\{\tau\hoch{ErbSt},\ 0{,}15\}\bigr],
  \qquad R\hoch{BoSt} = L\,\tau\hoch{BoSt},\\
  R\hoch{VSt} &= W\,\tau\hoch{VSt}\exp\bigl(-\varepsilon_W\tau\hoch{VSt}\bigr),\\
  R\hoch{Min} &= W^{*}\bigl(1 - 0{,}15\min\{1,\ \tau\hoch{Min}/0{,}02\}\bigr)\max\bigl\{0,\ \tau\hoch{Min} - 0{,}003 - \tau\hoch{VSt}\bigr\}.
\end{align}
Die Semi-Elastizität $\varepsilon_W = 21{,}5$, also 0,215 je Prozentpunkt, ist
die Hälfte des Schweizer Befunds, der Umzüge zwischen Kantonen einschließt.
Die Ausweichreaktion der Mindeststeuer, 15\,\% bei einem Satz von 2\,\%,
folgt \textcite{jakobsen2024}.

\paragraph{\COz-Preis.} Die Emissionen folgen einem Basispfad $\bar E(j)$ ohne
zusätzliche Politik, linear zwischen 649~Mt im Jahr 2025, $-63\,\%$ (2030) und
$-80\,\%$ (2040) gegenüber 1990 \parencite{uba2025}. Den nationalen Preis $p$
tragen nur die bepreisten Emissionen $\bar E^{p}(j) = \tfrac{327}{649}\bar E(j)$;
auf sie wirkt die Semi-Elastizität $\varepsilon_E = -0{,}30$ je 100\,\euro/t:
\begin{equation}
\begin{split}
  E_t &= \bar E(j_t) - \bar E^{p}(j_t)\bigl[1 - \psi(p)\bigr],\qquad
  R\hoch{CO_2} = \bar E^{p}(j_t)\,\psi(p)\,p,\\
  \psi(p) &= \min\bigl\{1{,}1,\ \max\{0{,}4,\ 1 + \varepsilon_E\,(p - 55)/100\}\bigr\}.
\end{split}
\end{equation}
Ein Klimageld zahlt 70\,\% von $R\hoch{CO_2}$ je Haushalt gleich aus; im
Status quo gibt es keines. Die übrigen Verbrauch- und Verkehrsteuern (Energie,
Tabak, Kfz, Solidaritätszuschlag und weitere, zusammen 109~Mrd.~\euro) stehen
fest.

\paragraph{Sozialbeiträge.} Renten-, Kranken-, Arbeitslosen- und
Pflegeversicherung erheben die Sätze $\beta\hoch{RV}$, $\beta\hoch{KV}$ und
$\beta\hoch{AP}$ (Arbeitslosen- und Pflegeversicherung) auf das
Arbeitseinkommen bis zu den Beitragsbemessungsgrenzen $\bar A\hoch{RV}$ und
$\bar A\hoch{KV} = \bar A\hoch{RV}\cdot 69.750/101.400$. Der Gesamtbeitrag aus
Arbeitnehmer- und Arbeitgeberanteil ist
\begin{equation}
  \sigma_i(\theta) = \min\{\tilde A_i,\ \bar A\hoch{RV}\}\bigl(\beta\hoch{RV} + 0{,}42\,\beta\hoch{AP}\bigr)
                   + \min\{\tilde A_i,\ \bar A\hoch{KV}\}\bigl(\beta\hoch{KV} + 0{,}58\,\beta\hoch{AP}\bigr).
\end{equation}
Der Staat rechnet mit der sozialversicherungspflichtigen Lohnsumme
$\Lambda = 1.750$~Mrd.~\euro{}, kalibriert bei den Grenzen des Status quo. Rund
12\,\% der beitragspflichtigen Löhne liegen zwischen der Grenze und knapp ihrem
Doppelten; eine um $x$\,\% andere Grenze ändert die Lohnsumme deshalb um
$0{,}12\,x$\,\%, in beide Richtungen:
\begin{equation}
  R\hoch{SV} = \Lambda\Bigl(1 + 0{,}12\,\frac{\bar A\hoch{RV} - \bar A^{\mathrm{RV},0}}{\bar A^{\mathrm{RV},0}}\Bigr)
              \bigl(\beta\hoch{RV} + \beta\hoch{KV} + \beta\hoch{AP}\bigr).
\end{equation}
Aus den Haushaltstypen lässt sich die Lohnmasse über der Grenze nicht rechnen: Die
Grenze gilt je Person, die Daten sind Haushaltseinkommen, und
Mehrverdienerhaushalte lägen fälschlich darüber. Die Beitragssätze bestimmen nur
die Einnahmen; die Leistungen hängen am Rentenniveau (Abschnitt~\ref{sec:ausgaben}).

\subsection{Ausgaben, Transfers und Renten}
\label{sec:ausgaben}

Die Ausgaben beginnen beim Stand von 2025 in grober Gliederung
(Tabelle~\ref{tab:ausgaben} im Anhang). Bürger- und Kindergeld sind Regelsatz
$\bar q$ beziehungsweise Betrag $\bar k$ je Monat mal Zahl der Berechtigten
aus den Profilen,
\begin{equation}
  X\hoch{BG} = 12\,\bar q\sum_i N_i q_i,\qquad X\hoch{KG} = 12\,\bar k\sum_i N_i k_i ;
\end{equation}
Unterkunft und Mehrbedarfe (10,8~Mrd.~\euro) bleiben fest. Die Renten folgen
dem Sicherungsniveau $\nu$ (Status quo 48\,\%, §\,154 Abs.~3 SGB~VI):
\begin{equation}
  \Delta P_i = \rho_i B_i f_t\bigl(\nu/\nu^0 - 1\bigr),\qquad \Delta X\hoch{R} = \sum_i N_i\,\Delta P_i .
\end{equation}
Die Profile $\rho_i$ sind so skaliert, dass die Renten der Haushalte den
Zahlungen der gesetzlichen Rentenversicherung entsprechen, rund
362~Mrd.~\euro{} \parencite{bmas2025}. Der Demografiefaktor $f_t$ bildet die
wachsende Zahl der Rentner nach der 14.~koordinierten
Bevölkerungsvorausberechnung ab \parencite{destatis2019}, von 1,00 im Jahr
2025 auf 1,25 im Jahr 2041. Er macht jede Niveauänderung teurer und erhöht die
Rentenausgaben um $X\hoch{Dem}_t = 390\,(f_t - 1)$~Mrd.~\euro.

Geltendes Recht hebt die Verteidigungsausgaben von 2,4\,\% auf 3,5\,\% des BIP
(2029) und finanziert aus dem Sondervermögen Infrastruktur von 2026 bis 2036
jährlich $500/11$~Mrd.~\euro{} öffentliche Investitionen (Art.~143h GG); $v_t$
ist die Verteidigungsquote der Periode, $O_t$ die Ausgabe des Sondervermögens.
Der Klima- und Transformationsfonds gibt aus, was der \COz-Preis des Status
quo auf dem Basispfad einbringt, und schrumpft mit den Emissionen. Die
Erhebungskosten sind feste Anteile am Aufkommen je Steuerart, etwa 2,5\,\%
der Einkommensteuer und 8\,\% der Erbschaftsteuer; in den Saldo geht ihre
Änderung $\Omega(\theta) - \Omega(\theta^0)$ ein.

\subsection{Inzidenz und Verteilung}
\label{sec:inzidenz}

Das verfügbare Einkommen eines Haushaltstyps ist
\begin{equation}
  Y_i = \tilde B_i - T_i - \hat\sigma_i - U_i - \ell_i + \Gamma_i + \Delta P_i - J_i .
  \label{eq:netto}
\end{equation}
Vier Posten folgen einer Inzidenzannahme. \emph{Beiträge:} Formal zahlen
Arbeitnehmer und Arbeitgeber je die Hälfte. Langfristig tragen die
Beschäftigten in Kontinentaleuropa rund zwei Drittel auch des
Arbeitgeberanteils \parencite{melguizo2013}. Das Modell lässt das Niveau beim
Arbeitnehmeranteil und rechnet jede Änderung voll dem Haushalt zu: bei den Sätzen
auch den Arbeitgeberanteil, bei der Beitragsbemessungsgrenze die ganze Buchung des
Staates, verteilt nach der Lohnmasse, die je Haushalt zwischen alter und neuer
Grenze liegt; $\hat\sigma_i$ sind die so getragenen Beiträge. \emph{Umsatzsteuer:} Eine Satzänderung trifft die volle
Bemessungsgrundlage,
$U_i = C_i\bigl[\bar\tau(\theta) + (\varphi(1-\gamma) - 1)\bigl(\bar\tau(\theta) - \bar\tau(\theta^0)\bigr)\bigr]$
mit dem Durchschnittssatz
$\bar\tau = 0{,}7\,\tau\hoch{USt}/(1+\tau\hoch{USt}) + 0{,}3\,\tau\hoch{erm}/(1+\tau\hoch{erm})$.
\emph{\COz:} Das Bruttoaufkommen wird vollständig überwälzt,
$\ell_i = \omega_i R\hoch{CO_2}$.
\emph{Unternehmen und Vermögen:} Ihre Steuern tragen am Ende Beschäftigte,
Eigentümer und Vermögende; von der Gewerbesteuer tragen die Beschäftigten
rund die Hälfte \parencite{fuest2018}. Eine Änderung von Körperschaft- und
Gewerbesteuer, $\Delta R\hoch{U}$, verteilt das Modell je zur Hälfte nach
Arbeits- und Kapitaleinkommen, eine Änderung von Erbschaft-, Vermögen- und
Bodenwertsteuer, $\Delta R\hoch{V}$, nach dem Vermögen, weil die
Bodenwertsteuer nicht überwälzbar ist \parencite[Kap.~16]{mirrlees2011}:
\begin{equation}
  J_i = \Delta R\hoch{U}\Bigl(\frac{A_i}{2\sum_l N_l A_l} + \frac{K_i}{2\sum_l N_l K_l}\Bigr)
      + \Delta R\hoch{V}\,\frac{V_i}{\sum_l N_l V_l},
\end{equation}
beide in Größen von 2025; die Mindeststeuer trägt allein D10c,
$J_{12}$ enthält zusätzlich $\Delta R\hoch{Min}/N_{12}$. Die Transfers
$\Gamma_i = 12\,(\bar q\,q_i + \bar k\,k_i) + \Gamma^{F}_i$ enthalten den festen Teil
der Grundsicherung $\Gamma^{F}_i$ und gegebenenfalls das Klimageld. Summiert
über die Haushalte ergeben Transfers, Renten, \COz-Last, Inzidenz und die Wirkung
der Beitragsbemessungsgrenze genau die Buchung des Staates, Umsatzsteuer und
Beitragssätze näherungsweise (Abschnitt~\ref{sec:validierung}).

Die Verteilungsmaße beziehen sich wie in EU-SILC auf das Äquivalenzeinkommen
$y_i = Y_i/w_i$ \parencite{oecd2013}. Der Gini-Koeffizient ist die mit $N_i$
gewichtete Fläche zwischen Lorenzkurve und Gleichverteilung, die
Palma-Ratio das Verhältnis der Einkommenssummen des obersten Dezils und der
unteren vier \parencite{palma2011}. Für die Armutsrisikoquote, den Anteil der
Haushalte unter 60\,\% des Medians, sind Durchschnitte zu grob: Kein
Haushaltstyp liegt im Mittel unter der Schwelle. Das Modell schreibt deshalb
die Anteile armutsgefährdeter Haushalte in D1 bis D3 aus dem Status quo fort,
\begin{equation}
  \pi_i = \min\Bigl\{1,\ \pi_i^0\Bigl(\frac{Y_i/Y_i^0}{Y^{*}/Y^{*0}}\Bigr)^{-1{,}5}\Bigr\},
  \qquad \pi^0 = (0{,}90;\ 0{,}62;\ 0{,}05),
\end{equation}
mit der Armutsschwelle $Y^{*}$, 60\,\% des mit $N_i$ gewichteten Medians der
$Y_i$, und der Elastizität $-1{,}5$ nach \textcite{bourguignon2003}.

\subsection{Nachfrage}
\label{sec:nachfrage}

Haushalte, die mehr verfügbares Einkommen haben als im Status quo derselben
Periode, geben einen Teil davon aus, unten mehr als oben. Zusammen mit
zusätzlichen öffentlichen Investitionen $\Delta O$ ergibt das den
Nachfrageimpuls $\Delta Z$ und die Nachfragelücke $u$ der Periode,
\begin{equation}
  \Delta Z = \frac{M_T}{\bar m}\sum_i N_i\,m_i\,\Delta Y_i + M_I\,\Delta O,
  \qquad u = \frac{\Delta Z}{\mathit{BIP}_0},
\end{equation}
mit dem Steuer- und Transfermultiplikator $M_T = 0{,}6$ bei mittlerer
Grenzkonsumneigung $\bar m = 0{,}48$, dem Investitionsmultiplikator
$M_I = 1{,}0$ \parencite{gechert2015} und dem BIP des Jahres 2025,
$\mathit{BIP}_0 = 4.470$~Mrd.~\euro. Die Lücke wirkt nur in der laufenden
Periode und symmetrisch: Konsolidierung kostet Wachstum, Entlastung bringt
welches \parencite{blanchard2013}. Von der Nachfrage bleibt dauerhaft nur,
was in öffentliches Kapital fließt.

\subsection{Saldo und Schuldenbremse}
\label{sec:saldo}

Bis hierher ist alles in Preisen von 2025 gerechnet. Die Einnahmen folgen dem
BIP der Periode mit einer Aufkommenselastizität von eins, einschließlich der
Nachfragelücke (automatischer Stabilisator); der Tarif gilt als indexiert.
Der \COz-Preis ist ein Preis und folgt dem Trend
$\Pi_t = (1+g)^{\,j_t - 2025}$ mit dem nominalen Wachstum $g = 2{,}5\,\%$. Die
Ausgaben wachsen mit dem Trend, nicht mit dem BIP; eine Rezession senkt sie
nicht. Zinsen entstehen aus dem Schuldenstand $D_t = d_t\,\mathit{BIP}_t$ zum
Effektivzins $z = 2{,}0\,\%$:
\begin{equation}
  R_t = \frac{\mathit{BIP}_t}{\mathit{BIP}_0}(1+u)\sum_{k \neq \mathrm{CO_2}} R^{k} + \Pi_t R\hoch{CO_2}_{\mathrm{netto}},
  \qquad
  X_t = \Pi_t\bigl(X^{\mathrm{real}} + V_t\bigr) + O_t + z\,D_t,
  \qquad
  S_t = R_t - X_t .
\end{equation}
$X^{\mathrm{real}}$ umfasst die Ausgaben ohne Zinsen in Preisen von 2025
einschließlich $X\hoch{BG}$, $X\hoch{KG}$, $\Delta X\hoch{R}$,
$X\hoch{Dem}_t$, $\Delta O$ und der Änderung der Erhebungskosten; $V_t$ ist
der Anstieg der Verteidigungsausgaben über 2,4\,\% des BIP, $R\hoch{CO_2}_{\mathrm{netto}}$
das \COz-Aufkommen nach einem Klimageld. Der strukturelle Saldo bereinigt um
die Konjunktur, mit der Budget-Semielastizität 0,5 der EU-Kommission
\parencite{mourre2019}, und um die Ausnahmen der Schuldenbremse von 2025,
Verteidigung über 1\,\% des BIP und das Sondervermögen:
\begin{equation}
  s^{\mathrm{str}}_t = \frac{S_t + O_t}{\mathit{BIP}_t(1+u)} - 0{,}5\,u + \max\{0,\ v_t - 0{,}01\}
  \;\ge\; -0{,}70\,\% .
\end{equation}
Das Modell rechnet den Gesamtstaat, nicht Bund und Länder getrennt.

\subsection{Von Periode zu Periode}
\label{sec:dynamik}

Das nominale BIP wächst mit dem Trend und drei Niveaueffekten der
Angebotsseite,
\begin{equation}
  \mathit{BIP}_{t+1} = \mathit{BIP}_t\,(1+g)^n\,
  \frac{\mathcal K_{t+1}}{\mathcal K_t}\,\frac{\mathcal A_{t+1}}{\mathcal A_t}\,\frac{\Phi(\mathcal Q_{t+1})}{\Phi(\mathcal Q_t)} .
\end{equation}
\emph{Arbeit} wirkt ab der folgenden Periode ohne Anpassungsverzögerung, mit dem
Arbeitsanteil 0,65,
$\mathcal A_{t+1} = 1 + 0{,}65\,(\bar\lambda_t - 1)$, wobei $\bar\lambda_t$ das
mit $N_i$ gewichtete Mittel der $\lambda_i$ ist. \emph{Privates Kapital}
nähert sich mit $\delta = 7\,\%$ je Jahr dem Niveau, das die
Unternehmensteuern langfristig tragen,
$\mathcal K_{t+1} = \mathcal K_t + [1 - (1-\delta)^n]\,(\mathcal K^{*}_t - \mathcal K_t)$
mit $\mathcal K^{*}_t = 1 + 0{,}35\,(I_t - 1)$. \emph{Öffentliches Kapital}
$\mathcal Q$ aus zusätzlichen Investitionen und dem Sondervermögen schreibt mit
4\,\% je Jahr ab und hebt das Potenzial mit der Produktionselastizität 0,08,
$\Phi(\mathcal Q) = 1 + 0{,}08\,\mathcal Q/1.600$~Mrd.~\euro{}
\parencite{bom2014}. Alle drei sind Niveaus, keine Wachstumsraten: Eine höhere
Körperschaftsteuer senkt das BIP auf einen niedrigeren Pfad, aber nicht jede
Periode weiter.

Die Schulden wachsen Jahr für Jahr mit dem Zins und fallen mit dem
Primärsaldo der Periode, die Emissionen summieren sich entlang des
Basispfads:
\begin{equation}
  D \leftarrow (1+z)\,D - \bigl(S_t + z\,D_t\bigr)\ \ (n\text{-mal}),
  \qquad d_{t+1} = \frac{D_{t+1}}{\mathit{BIP}_{t+1}},
  \qquad \mathcal E_{t+1} = \mathcal E_t + \frac{E_t}{\bar E(j_t)}\sum_{j\in t}\bar E(j).
\end{equation}
Die Lohnsumme $\Lambda$ folgt dem Arbeitsangebot mit Rückkehr zum Trend.
Exogene Schocks verschieben BIP, Schuldenquote oder Zins.

\paragraph{Tragfähigkeit.} Weil der Zins unter der nominalen Wachstumsrate
liegt, $z - g = -0{,}5$~Pp., stabilisiert sich die Schuldenquote schon bei
einem kleinen Primärdefizit \parencite{domar1944,blanchard2019}. Das Modell
berichtet den Primärsaldo $ps_t = s_t + z\,d_t$, den
schuldenstabilisierenden Primärsaldo $ps^{*}_t = (z - g)\,d_t$ und die Lücke
$ps_t - ps^{*}_t$, dazu einen Index der Generationengerechtigkeit. Er bildet
die Schuldenquote zwischen 60 und 150\,\% und das verbrauchte \COz-Budget zu
gleichen Teilen auf $[0,1]$ ab; das deutsche Budget von 5.380~Mt ab 2025 ist
der Bevölkerungsanteil am globalen Budget für 1,7\,°C \parencite{forster2026}.
